\documentclass[aspectratio=169]{beamer}
% Shared Beamer style for this repository.
\usepackage[fontset=none]{ctex}
\usepackage{fontspec}
\setmainfont{DejaVu Serif}
\setsansfont{DejaVu Sans}
\setmonofont{DejaVu Sans Mono}
\setCJKmainfont{Noto Serif CJK SC}[BoldFont={Noto Serif CJK SC Bold}]
\setCJKsansfont{Noto Sans CJK SC}[BoldFont={Noto Sans CJK SC Bold}]
\setCJKmonofont{Noto Sans Mono CJK SC}

\usetheme{Madrid}
\usefonttheme{serif}
\setbeamertemplate{navigation symbols}{}
\setbeamertemplate{footline}[frame number]

\usepackage{graphicx}
\usepackage{amsmath,amssymb}
\usepackage{booktabs}
\usepackage{array}
\usepackage{xcolor}
\usepackage{listings}
\usepackage{tikz}
\usetikzlibrary{shapes,arrows,positioning}

% Color system
\definecolor{dcPrimary}{RGB}{21,101,117}
\definecolor{dcSecondary}{RGB}{140,115,72}
\definecolor{dcAccent}{RGB}{184,92,63}
\definecolor{dcMuted}{RGB}{230,250,252}
\definecolor{dcWarmBG}{RGB}{255,245,230}
\definecolor{dcCodeBG}{RGB}{247,248,250}
\definecolor{dcCodeFrame}{RGB}{200,210,215}

\setbeamercolor{structure}{fg=dcPrimary}
\setbeamercolor{palette primary}{fg=dcPrimary,bg=dcMuted}
\setbeamercolor{palette secondary}{fg=dcSecondary,bg=dcMuted!90}
\setbeamercolor{palette tertiary}{fg=dcPrimary,bg=dcMuted!80}
\setbeamercolor{frametitle}{fg=white,bg=dcPrimary}
\setbeamercolor{title}{fg=dcPrimary}
\setbeamercolor{alerted text}{fg=dcAccent}
\setbeamercolor{item}{fg=dcPrimary}
\setbeamercolor{subitem}{fg=dcSecondary}
\setbeamercolor{block title}{fg=white,bg=dcPrimary}
\setbeamercolor{block body}{bg=dcMuted}
\setbeamercolor{block title example}{fg=white,bg=dcSecondary}
\setbeamercolor{block body example}{bg=dcWarmBG}
\setbeamercolor{block title alerted}{fg=white,bg=dcAccent}
\setbeamercolor{block body alerted}{bg=dcWarmBG}
\setbeamercolor{footline}{fg=dcPrimary}

\lstset{
  basicstyle=\ttfamily\footnotesize,
  backgroundcolor=\color{dcCodeBG},
  frame=single,
  rulecolor=\color{dcCodeFrame},
  numbers=none,
  breaklines=true,
  showstringspaces=false,
  tabsize=2
}


\title{<Lecture Title>}
\subtitle{<Lecture Subtitle>}
\author{<Instructor>}
\date{\today}

\graphicspath{{figures/}}

\begin{document}

\begin{frame}
  \titlepage
\end{frame}

\begin{frame}{Learning Objectives}
% 建议 3--5 条，说明读者这一讲结束后会多出什么判断框架。
\begin{itemize}
  \item <Objective 1>
  \item <Objective 2>
  \item <Objective 3>
\end{itemize}
\end{frame}

\begin{frame}{Why This Lecture}
% 先说清：为什么有这一讲？它和前后章节是什么关系？
% PPT 也要先给地图，不要第一页正文就开始堆定义。
<Explain why this lecture exists in the course map.>
\end{frame}

\begin{frame}{Core Map or Key Concept}
% 一页一主点。优先放最关键的概念图、层次图、流程图或关系图。
% 如果这是第一次引入概念，先立轮廓，不要把所有细节塞进第一页。
<Main concept, map, or diagram here.>
\end{frame}

\begin{frame}{Example or System View}
% 用一个例子把抽象概念放回真实对象或真实系统中。
<Example, case study, or figure explanation here.>
\end{frame}

\begin{frame}{Common Confusion or Boundary}
% 这一页专门放容易混淆的点、边界条件或常见误解。
\begin{itemize}
  \item <Confusion 1>
  \item <Boundary 1>
  \item <Misconception 1>
\end{itemize}
\end{frame}

\begin{frame}{Takeaways}
% 回收本讲主线，不要只是重复目录。
\begin{itemize}
  \item <Takeaway 1>
  \item <Takeaway 2>
  \item <Takeaway 3>
\end{itemize}
\end{frame}

\begin{frame}{After-Class Prompt}
% 课后提示优先设计成“画图 / 分类 / 解释 / 对比 / 分析”任务。
<One short after-class prompt here.>
\end{frame}

\end{document}
